\begin{tikzpicture}[scale=0.8,font=\small]
    % === LÍNEAS DE REFERENCIA ===
    % Ejes y límites de las áreas asintóticas
    \draw[red!0, name path=A--A'] (0,6) -- ++(6,0);
    \draw[red!0, name path=C--C'] (0,8) -- ++(11,0);

    % === DEFINICIÓN DE LA CURVA DE ESFUERZO ===
    % Trazado paramétrico de la curva principal
    \path[name path=f--f'] (0,0) .. controls (1,3) and (3,8) .. (11,8);

    % --- Cálculo de Intersecciones ---
    \path[name intersections={of=A--A' and f--f', by=B}];
    \path[name intersections={of=C--C' and f--f', by=D}];

    % === DEFINICIÓN DE LA CURVA DE RELAJACIÓN ===
    % Trazado secundario de descarga
    \path[name path=f2--f2'] (1.5,0) .. controls (3,5) and (4.35,8) .. (D);

    % --- Cálculo de Intersección Central ---
    \path[name intersections={of=f--f' and f2--f2',by=E}];

    % === SOMBREADO DE ÁREAS (USANDO CLIP) ===
    % 1. Área elástica (Azul)
    \begin{scope}
        \clip (0,0) rectangle (B);
        \fill[azul!20] (0,0) .. controls (1,3) and (3,8) .. (11,8) -- (0,6) -- cycle;
    \end{scope}

    % 2. Área plástica (Verde con opacidad)
    \begin{scope}
        \clip (0,6) rectangle (13,8);
        \fill[verde!40,opacity=0.7] (0,0) .. controls (1,3) and (3,8) .. (11,8) -- (0,8) -- cycle;
    \end{scope}

    % 3. Área de histéresis/relajación
    \begin{scope}
        \clip (0,0) .. controls (1,3) and (3,8) .. (11,8) -- (11,0) -- (0,0) -- cycle;
        \draw[curva vectorial=grisa] (E) .. controls (4,6) and (3,4) .. (1.5,0);
    \end{scope}

    % === DIBUJO DE LAS CURVAS (CON DECORATIONS) ===
    \draw[curva vectorial=mainc] (0,0) .. controls (1,3) and (3,8) .. (D);

    % --- Puntos Singulares ---
    \fill[yellow] (B) circle (0.7mm); % Límite elástico
    \fill[rojo] (D) circle (0.7mm);   % Ruptura

    % === MARCAS DE LOS EJES (TICKS) ===
    % Eje Y
    \foreach \y in {0.5,1,...,8} {
        \draw[shift={(0,\y)},color=black] (0,0) -- (-1.1mm,0);
    }
    % Eje X
    \foreach \x in {0.5,1,...,11.5} {
        \draw[shift={(\x,0)},color=black] (0,0) -- (0,-1.1mm);
    }

    % === DIBUJO DE LOS EJES PRINCIPALES ===
    \draw[-stealth,very thick] (-0.5,0) -- ++(12.5,0) node[below] {$Du$ $(\times 10^{-4})$};
    \draw[-stealth,very thick] (0,-0.5) -- ++(0,9) node[above] {$E$ (MPa)};

    % === LEYENDA DEL GRÁFICO ===
    \matrix[
        draw=gray,
        line width=0.4mm,
        fill=white,
        right,
        at={($(current bounding box.center)+(2,1)$)},
        nodes={anchor=west},
        inner sep=1mm,
        column sep=1mm,
        rounded corners,
        column 1/.style={nodes={anchor=center}}
    ]
    {
        % Fila 1: Curva principal
        \draw[vector,mainc,ultra thick] (-0.2,0) -- (0.2,0); &
        \node {Curva de Esfuerzo}; \\

        % Fila 2: Curva secundaria
        \draw[vector,grisa,ultra thick] (-0.2,0) -- (0.2,0); &
        \node {Curva de Relajación}; \\

        % Fila 3: Punto de transición
        \fill[yellow] (0,0) circle (0.7mm); &
        \node {Límite Elástico}; \\

        % Fila 4: Punto de fallo
        \fill[rojo] (0,0) circle (0.7mm); &
        \node {Ruptura del material}; \\
    };
\end{tikzpicture}